\documentclass[11pt,letterpaper]{article}
\usepackage[margin=1in]{geometry}
\usepackage{enumitem}
\usepackage{hyperref}
\usepackage{amsmath,amssymb}
\usepackage{xcolor}

\hypersetup{colorlinks=true, linkcolor=blue!60!black, urlcolor=blue!60!black}

\newenvironment{keypoint}{\medskip\noindent\fcolorbox{blue!40!black}{blue!5!white}{\begin{minipage}{\dimexpr\textwidth-2\fboxsep-2\fboxrule}\textbf{Key Message:} }{\end{minipage}}\medskip}
\newenvironment{figureidea}{\medskip\noindent\fcolorbox{green!40!black}{green!5!white}{\begin{minipage}{\dimexpr\textwidth-2\fboxsep-2\fboxrule}\textbf{Figure Idea:} }{\end{minipage}}\medskip}

\title{\textbf{Paper Outline: Rayfield} \\ \large Physics-Embedded Flow Matching for Photonic Device Field Prediction}
\author{Paper Structure Guide}
\date{March 2026}

\begin{document}
\maketitle
\tableofcontents
\newpage

% ===========================================================================
\section{1. Title \& Abstract}
% ===========================================================================

\subsection{Title Options}
\begin{itemize}
    \item ``Rayfield: Physics-Embedded Flow Matching for Electromagnetic Field Prediction in Photonic Integrated Circuits''
    \item ``Complex-Valued Flow Matching with Helmholtz Constraints for Photonic Device Simulation''
    \item Shorter: ``Physics-Constrained Generative Modeling for Photonic Field Prediction''
\end{itemize}

\subsection{Abstract ($\sim$150--250 words)}
\begin{enumerate}
    \item \textbf{Problem}: FDTD simulation of photonic devices is slow (minutes--hours per geometry), limiting design-space exploration and inverse design.
    \item \textbf{Gap}: Existing surrogate models either ignore physics (pure data-driven) or are limited to simple geometries. No prior work combines flow matching + complex-valued networks + Helmholtz constraints for photonic simulation.
    \item \textbf{Method}: We present a complex-valued Physics UNet trained via conditional flow matching with a multi-loss curriculum (FM $\to$ Helmholtz residual $\to$ S-parameters). Native complex arithmetic preserves phase. Helmholtz residual loss enforces Maxwell's equations.
    \item \textbf{Results}: Quote key metrics --- PSNR, SSIM, S-parameter accuracy, and speedup vs.\ FDTD. Mention number of device types and dataset size.
    \item \textbf{Impact}: Enables real-time field prediction for inverse design of photonic circuits.
\end{enumerate}

% ===========================================================================
\section{2. Introduction}
% ===========================================================================

\subsection{2.1 Motivation \& Context ($\sim$1 page)}
\begin{itemize}
    \item Photonic integrated circuits (PICs) are critical for telecom, sensing, quantum computing, AI accelerators.
    \item Design relies on electromagnetic simulation (FDTD, FEM) --- accurate but computationally expensive.
    \item Growing need for fast surrogate models to enable: (a) large-scale design-space exploration, (b) gradient-based inverse design, (c) real-time design feedback.
\end{itemize}

\subsection{2.2 Limitations of Prior Work}
\begin{itemize}
    \item \textbf{Physics-informed neural networks (PINNs)}: Require per-geometry optimization; don't generalize across geometries.
    \item \textbf{Data-driven surrogates}: CNNs/U-Nets that predict fields directly --- ignore physics, poor generalization, phase errors.
    \item \textbf{Neural operators} (FNO, DeepONet): Promising but typically real-valued; lose phase information critical for coherent photonics.
    \item \textbf{Generative models}: Diffusion/flow matching used in other domains but not yet applied to EM field prediction with physics constraints.
\end{itemize}

\subsection{2.3 Contributions}
\begin{enumerate}
    \item \textbf{Complex-valued flow matching}: First application of conditional flow matching with native complex arithmetic to EM field prediction.
    \item \textbf{Helmholtz-constrained training}: Physics residual loss that enforces $\nabla^2 E + k_0^2 \varepsilon E = 0$ with interface-aware masking.
    \item \textbf{Multi-loss curriculum}: Phased training (FM $\to$ physics $\to$ S-params) that stabilizes convergence.
    \item \textbf{Multi-device generalization}: Single model trained across 5+ photonic device types with varying port counts and physics.
    \item \textbf{Dataset}: Large-scale multi-device FDTD dataset (10,000 geometries $\times$ 4 wavelengths = 40,000 simulations) with standardized format.
\end{enumerate}

% ===========================================================================
\section{3. Background}
% ===========================================================================

\subsection{3.1 Electromagnetic Simulation of Photonic Devices}
\begin{itemize}
    \item Maxwell's equations in 2D TE polarization $\to$ scalar Helmholtz: $\nabla^2 E_z + k_0^2 \varepsilon(x,y) E_z = 0$.
    \item FDTD method: time-stepping on Yee grid, converges to frequency-domain solution.
    \item Key outputs: complex field $E_z(x,y)$, permittivity map $\varepsilon(x,y)$, S-parameters.
    \item Silicon photonics material system: Si core ($n_\text{eff} \approx 2.4$), SiO$_2$ cladding ($n = 1.444$), C-band wavelengths (1500--1600\,nm).
\end{itemize}

\subsection{3.2 Conditional Flow Matching}
\begin{itemize}
    \item Flow matching as an alternative to diffusion: learns a velocity field $u_t(x_t, t)$ that transports noise $x_0 \sim \mathcal{N}(0,I)$ to data $x_1$.
    \item Interpolation path: $x_t = (1 - (1-\sigma)t)\,x_0 + t\,x_1$.
    \item Training objective: $\mathcal{L}_\text{FM} = \mathbb{E}_{t,x_0,x_1}\left[\|u_\theta(x_t, t) - u_t\|^2\right]$.
    \item Advantages: simpler training than diffusion, deterministic sampling via ODE integration.
    \item Cite: Lipman et al.\ (2023), Liu et al.\ (2023).
\end{itemize}

\subsection{3.3 Physics-Informed Neural Networks}
\begin{itemize}
    \item Brief survey of PINNs, physics-constrained learning.
    \item Distinction: PINNs optimize per-instance; our approach trains a generalizable model with physics loss as regularization.
    \item Cite: Raissi et al.\ (2019), and relevant photonics PINN papers.
\end{itemize}

% ===========================================================================
\section{4. Method}
% ===========================================================================

\begin{keypoint}
This is the core of the paper. Expect 4--5 pages with equations and architecture diagrams.
\end{keypoint}

\subsection{4.1 Problem Formulation}
\begin{itemize}
    \item \textbf{Input}: Permittivity map $\varepsilon(x,y)$ [1,H,W], source mask [1,H,W], wavelength $\lambda$.
    \item \textbf{Output}: Complex field $E_z(x,y) = E_r + jE_i$ [2,H,W].
    \item \textbf{Auxiliary targets}: S-parameters (complex-valued, per-port), port masks.
    \item Conditioning: wavelength (normalized scalar) + geometric design parameters.
\end{itemize}

\subsection{4.2 Complex-Valued Physics UNet}
\begin{itemize}
    \item \textbf{Architecture}: U-Net with encoder--bottleneck--decoder + skip connections.
    \item \textbf{Complex layers}:
    \begin{itemize}
        \item \texttt{ComplexConv2d}: $(W_r + jW_i)(x_r + jx_i) = (W_r x_r - W_i x_i) + j(W_r x_i + W_i x_r)$
        \item \texttt{ComplexGroupNorm}: Joint normalization of real/imaginary parts.
        \item \texttt{ModReLU}: $f(z) = \text{ReLU}(|z| + b) \cdot \frac{z}{|z|}$ --- gates magnitude, preserves phase.
        \item \texttt{ComplexAttention}: Q, K, V in complex domain.
    \end{itemize}
    \item \textbf{Timestep conditioning}: Sinusoidal embedding + MLP, injected via AdaGN (adaptive group norm).
    \item \textbf{Physics feature embedding}: Helmholtz residual features computed inside the network as additional input channels.
    \item Discuss why complex-valued: phase is first-class citizen, avoids phase wrapping issues in real-valued models.
\end{itemize}

\begin{figureidea}
Architecture diagram showing the Complex UNet with: (1) input channels, (2) encoder/decoder blocks with complex layers, (3) skip connections, (4) timestep injection, (5) physics feature computation.
\end{figureidea}

\subsection{4.3 Flow Matching Formulation}
\begin{itemize}
    \item Conditional flow matching on the field channels [Re($E_z$), Im($E_z$)] only.
    \item Geometry ($\varepsilon$, source mask) as conditioning --- not noised.
    \item Interpolation: $x_t = (1-(1-\sigma_\text{min})t)\,x_0 + t\,x_1$, with $\sigma_\text{min} = 10^{-4}$.
    \item Time sampling: 50/50 mixture of uniform and Beta(2,2) $\to$ concentrates training near $t=0.5$.
    \item ODE sampling at inference: Euler or adaptive solver over $t \in [0,1]$.
\end{itemize}

\subsection{4.4 Physics-Constrained Losses}

\subsubsection{4.4.1 Helmholtz Residual Loss}
\begin{itemize}
    \item Residual: $R = \nabla^2 E_z + k_0^2 \varepsilon E_z$ computed via 5-point finite-difference Laplacian.
    \item Applied to the model's predicted $\hat{x}_1$ (endpoint estimate), not the noisy $x_t$.
    \item \textbf{Spatial masking}: Exclude PML regions, source regions, and \textbf{dielectric interface pixels} (where $|\nabla\varepsilon| > \tau$).
    \item \textbf{Interface masking rationale}: The collocated-grid Laplacian stencil produces $O(1)$ errors at sharp $\varepsilon$ discontinuities (Si/SiO$_2$ boundaries). Masking these pixels drops the GT residual floor from $\sim$0.17 to $\sim$0.01--0.03, giving a meaningful physics metric.
    \item \textbf{Normalization}: Divide by driving-term scale $\langle(k_0^2 \varepsilon E)^2\rangle$ to make the loss dimensionless and grid-independent.
    \item \textbf{Time gating}: Only apply when $t \geq t_\text{min}$ (e.g., 0.5) since low-$t$ predictions are too noisy for meaningful physics evaluation.
    \item Optional $t$-power weighting: upweight high-$t$ samples where field estimates are cleaner.
\end{itemize}

\subsubsection{4.4.2 S-Parameter Loss}
\begin{itemize}
    \item S-parameters extracted via modal decomposition: project predicted field onto waveguide eigenmodes at each port.
    \item Loss on both magnitude and phase of S-parameters vs.\ FDTD ground truth.
    \item Discuss port mask conventions and multi-port handling (2-port, 3-port, 4-port devices).
\end{itemize}

\subsubsection{4.4.3 Endpoint Loss}
\begin{itemize}
    \item Direct MSE between predicted and true fields at $t=1$.
    \item Amplitude-weighted to focus on high-signal regions.
\end{itemize}

\subsection{4.5 Multi-Loss Curriculum Training}
\begin{itemize}
    \item \textbf{Phase A} (epochs 0--$N_A$): FM loss only. Model learns basic field structure and velocity prediction.
    \item \textbf{Phase B} (epochs $N_A$--$N_B$): Add Helmholtz residual + endpoint loss. Physics constraints guide the model toward valid solutions.
    \item \textbf{Phase C} (epochs $N_B$+): Add S-parameter loss. Fine-tune port coupling accuracy.
    \item Each loss has independent warmup ramp to prevent gradient shock.
    \item Rationale: Introducing all losses simultaneously causes gradient conflict; curriculum stabilizes training.
\end{itemize}

\begin{figureidea}
Training curriculum timeline showing which losses are active at each phase, with warmup ramps. Could also show loss curves over training.
\end{figureidea}

\subsection{4.6 Data Augmentation}
\begin{itemize}
    \item D2 symmetry group (horizontal flip) applied on-the-fly during training.
    \item Exploits the vertical mirror symmetry of elongated photonic devices.
    \item Effectively doubles training data without additional FDTD computation.
    \item Note: Full D4 (rotations) not applicable due to rectangular grid and non-square domain.
\end{itemize}

% ===========================================================================
\section{5. Dataset}
% ===========================================================================

\subsection{5.1 Device Types}
\begin{itemize}
    \item Table of all device types with: name, number of ports, key parameters, parameter ranges.
    \item Devices: S-bend, Y-branch, directional coupler (+ optionally straight waveguide, MMI, euler bend, circular bend, crossing).
    \item Silicon-on-insulator platform: $n_\text{eff} \approx 2.4$ (core), $n_\text{clad} = 1.444$, C-band.
\end{itemize}

\subsection{5.2 FDTD Simulation Setup}
\begin{itemize}
    \item Meep FDTD solver, 2D TE polarization ($E_z$ component).
    \item Resolution: pixels/$\mu$m (discuss tradeoff with data size).
    \item PML absorbing boundaries.
    \item Eigenmode source excitation.
    \item Convergence criterion: field decay tolerance.
    \item Per-sample outputs: $\varepsilon(x,y)$, complex $E_z(x,y)$, source mask, port masks, S-parameters.
\end{itemize}

\subsection{5.3 Dataset Statistics}
\begin{itemize}
    \item Total geometries, wavelengths per geometry, total simulations.
    \item Latin hypercube sampling of parameter space with half-pixel quantization.
    \item Train/val/test split by geometry ID (80/10/10) --- no wavelength leakage.
    \item Data format: sharded NPZ files with JSON index.
\end{itemize}

\begin{figureidea}
(1) Example device geometries (one per type) showing $\varepsilon$ maps. (2) Corresponding $|E_z|$ fields. (3) Parameter distribution histograms.
\end{figureidea}

% ===========================================================================
\section{6. Experiments}
% ===========================================================================

\subsection{6.1 Training Setup}
\begin{itemize}
    \item Hardware: GPU type, number of GPUs, DDP setup.
    \item Optimizer: AdamW, learning rate schedule (warmup + cosine decay).
    \item Mixed precision (FP16 forward, FP32 losses).
    \item Hyperparameters: hidden size, channel multipliers, attention resolutions, batch size, epochs.
    \item Curriculum timing: Phase A/B/C epoch boundaries and loss weights.
\end{itemize}

\subsection{6.2 Baselines}
\begin{itemize}
    \item \textbf{Real-valued UNet}: Same architecture but with standard real convolutions (ablation of complex layers).
    \item \textbf{No physics loss}: FM-only training (ablation of Helmholtz constraint).
    \item \textbf{No curriculum}: All losses from epoch 0 (ablation of phased training).
    \item \textbf{Direct regression}: UNet trained with MSE loss only (no flow matching).
    \item Optionally: compare against FNO or other neural operator baselines.
\end{itemize}

\subsection{6.3 Evaluation Metrics}
\begin{itemize}
    \item \textbf{Field quality}: PSNR, SSIM, relative $L_2$ error (amplitude and phase separately).
    \item \textbf{Physics compliance}: Helmholtz residual ratio (with interface masking) --- report GT floor and model value.
    \item \textbf{S-parameter accuracy}: Magnitude error (dB), phase error (degrees), per-port breakdown.
    \item \textbf{Generalization}: Per-device-type metrics, out-of-distribution performance.
    \item \textbf{Computational cost}: Inference time vs.\ FDTD, model parameter count, FLOPs.
\end{itemize}

\subsection{6.4 Results}

\subsubsection{6.4.1 Field Prediction Quality}
\begin{itemize}
    \item Quantitative: Table of PSNR, SSIM, $L_2$ error per device type.
    \item Qualitative: Side-by-side comparison of GT vs.\ predicted fields (amplitude and phase).
    \item Show that complex UNet preserves phase structure better than real-valued baseline.
\end{itemize}

\begin{figureidea}
Grid of example predictions: rows = device types, columns = $\varepsilon$, $|E_z|$ GT, $|E_z|$ pred, phase GT, phase pred, error map.
\end{figureidea}

\subsubsection{6.4.2 Physics Compliance}
\begin{itemize}
    \item Helmholtz residual: model vs.\ GT floor, with and without interface masking.
    \item Show that physics loss reduces residual beyond what FM-only achieves.
    \item Spatial maps of residual --- should be near-zero away from interfaces.
\end{itemize}

\subsubsection{6.4.3 S-Parameter Accuracy}
\begin{itemize}
    \item Scatter plots: predicted vs.\ true $|S_{ij}|$ and $\angle S_{ij}$.
    \item Per-device breakdown (especially directional coupler which has strong wavelength dependence).
    \item Compare S-param accuracy with and without S-param loss.
\end{itemize}

\subsubsection{6.4.4 Ablation Studies}
\begin{itemize}
    \item Complex vs.\ real-valued UNet.
    \item With vs.\ without Helmholtz residual loss.
    \item With vs.\ without curriculum (all losses from start vs.\ phased).
    \item Effect of interface masking on residual metric.
    \item Effect of time gating ($t_\text{min}$) on physics loss effectiveness.
    \item Effect of $t$-power weighting on residual convergence.
\end{itemize}

\subsubsection{6.4.5 Computational Efficiency}
\begin{itemize}
    \item Inference time: single sample, batched.
    \item Speedup factor vs.\ FDTD (expect 100--1000$\times$).
    \item Training cost: GPU-hours, wall-clock time.
\end{itemize}

% ===========================================================================
\section{7. Discussion}
% ===========================================================================

\begin{itemize}
    \item \textbf{Why complex-valued matters}: Phase preservation for coherent devices (interferometers, couplers). Real-valued models struggle with phase wrapping.
    \item \textbf{Physics loss as regularization}: The Helmholtz loss doesn't just improve metrics --- it constrains the solution space to physically valid fields, improving generalization.
    \item \textbf{Interface masking insight}: Finite-difference Helmholtz residual has an inherent floor at dielectric interfaces due to stencil errors. Masking these pixels ($|\nabla\varepsilon| > \tau$) gives a meaningful metric. Discuss implications for any work using PDE residuals on discontinuous-coefficient problems.
    \item \textbf{Curriculum necessity}: Multi-task gradient conflicts between FM velocity loss and physics losses. Phased introduction prevents early-training instability.
    \item \textbf{Limitations}:
    \begin{itemize}
        \item 2D only (TE polarization) --- extension to 3D is future work.
        \item Fixed grid resolution --- does not adapt to feature size.
        \item Single-frequency per sample (not broadband).
        \item Trained on specific device families --- generalization to arbitrary freeform geometries untested.
    \end{itemize}
    \item \textbf{Broader impact}: Fast field prediction enables gradient-based inverse design, topology optimization, and automated PIC layout.
\end{itemize}

% ===========================================================================
\section{8. Related Work}
% ===========================================================================

\begin{itemize}
    \item \textbf{Neural surrogates for EM simulation}: Review of CNN/UNet-based field predictors (cite relevant photonics ML papers).
    \item \textbf{Physics-informed learning for photonics}: PINNs applied to waveguides, resonators, etc.
    \item \textbf{Neural operators}: FNO, DeepONet applied to PDEs; any photonics applications.
    \item \textbf{Flow matching / diffusion for scientific computing}: Score-based models for PDE solutions, molecular dynamics, weather prediction.
    \item \textbf{Complex-valued neural networks}: History and recent applications in signal processing, MRI reconstruction, wave physics.
    \item \textbf{Inverse design of photonic devices}: Topology optimization, adjoint methods, and how surrogates could accelerate these.
\end{itemize}

% ===========================================================================
\section{9. Conclusion}
% ===========================================================================

\begin{itemize}
    \item Summarize: complex-valued flow matching + Helmholtz constraint + curriculum training = fast, physics-compliant photonic field prediction.
    \item Restate key numbers: accuracy, speedup, number of device types.
    \item \textbf{Future work}:
    \begin{itemize}
        \item 3D simulations (full vectorial fields).
        \item Broadband prediction (multiple wavelengths in one forward pass).
        \item Integration with inverse design / topology optimization loops.
        \item Freeform geometry generalization (not just parameterized devices).
        \item Larger models / more data scaling laws.
    \end{itemize}
\end{itemize}

% ===========================================================================
\section{10. Appendix (Supplementary Material)}
% ===========================================================================

\begin{itemize}
    \item \textbf{A. Full architecture details}: Layer counts, channel dimensions, attention configuration.
    \item \textbf{B. Device parameter ranges}: Complete table of all parameter ranges per device type.
    \item \textbf{C. Training hyperparameters}: Full table of all hyperparameters used.
    \item \textbf{D. Additional qualitative results}: More example predictions across device types and wavelengths.
    \item \textbf{E. Phase anchoring}: Details of the global phase alignment procedure used before computing phase metrics.
    \item \textbf{F. Modal S-parameter extraction}: Mathematical details of eigenmode projection for S-parameter computation.
    \item \textbf{G. Helmholtz residual derivation}: Units, interface masking threshold selection, GT floor analysis.
\end{itemize}

\vfill
\begin{center}
\textit{Generated for the Rayfield project --- March 2026}
\end{center}

\end{document}
